% \documentclass[a4paper]{article}
% \usepackage{ctex}
\documentclass[AutoFakeBold]{ZafuThesis}

\usepackage[sort]{gbt7714}%参考文献
\setmainfont{Times New Roman}
\setCJKmainfont{宋体}

\title{浙江农林大学本科生毕业设计 说明书(论文)模板}
\academy{风景园林与建筑学院}
\session{2024}
\author{毛神神}
\studentNumber{201900000000}
\myClass{建筑学191}
\teacher{天神}{讲师}
\date{\today}
\title{{浙江农林大学本科生毕业设计 }{说明书(论文)模板}}
\entitle{Architectural design of Heshang Police Station of Xiaoshan District Public Security Bureau}


\begin{document}

% %封面
\customCover

% 承诺页
\makestatement
%自己的签名
\mysignature{\includegraphics[width=100pt]{zafu}}

%目录
\customContent

\frontmatter
%中文摘要，关键词分号分隔
\ZhAbstract
{
  摘要内容主要介绍所研究的课题内容、提出主要结论及创新之处。摘要部分格式：黑体加粗五号段前空两个汉字字符，摘要内容楷体五号，不超过300个汉字。关键词部分格式：黑体加粗五号段前空两个汉字字符，关键词内容楷体五号，术语用分号隔开，数量一般为3-6个。Abstract：Times New Roman加粗五号段前空两个汉字字符；Abstract内容：Times New Roman五号。Key Words: Times New Roman加粗五号段前空两个汉字字符； Key Words内容：Times New Roman五号，术语用逗号隔开。
  
}
{关键词1；关键词2；关键词3；关键词4；关键词5}


%英文摘要，关键词用逗号分隔
\EnAbstract
{In this short article we will discuss about LATEX for your dissertation}
{keyword 1, keyword 2, keyword 3}



\mainmatter
\section{模版介绍与注意事项}
\subsection{模版说明}
ZafuTemplate是作者本人自己写毕业论文时心血来潮所作，现发布到网上供浙江农林大学本科毕业生免费使用，在使用前请仔细阅读下面的内容和注意事项。\par
注意：本模版尚存许多不足之处，欢迎大家反馈，更希望大家能帮忙一起完善这个模版。
\subsection{下载安装}
ZafuTemplate 主页:...。
\subsection{文件说明}
在/ZafuTemplate目录中含有以下几个主要的文件，使用者可以根据自己的需要进行修改，从而完成自己的毕业论文：
\begin{itemize}
  \item these.tex：论文的内容，撰写论文请在该文件下。
  \item ZafuThesis.cls：用来控制论文格式的文件，如果有需要，可以自行修改格式。
\end{itemize}

\subsection{参考文献}
\subsection{编译注意事项}
\subsection{免责声明}
本模板依据《浙江农林大学本科生毕业论文（设计）系列材料（2022版）》编写，作者希 望能给使用者写作论文带来方便。然而，作者不保证本模板完全符合学校要求，也 不对由此带来的风险和损失承担任何责任。
\section{格式要求}
\subsection{论文格式参考}
本部分参阅：《浙江农林大学本科生毕业论文（设计）系列材料（2022版）》，以下将罗列出学校论文具体格式要求，供模版使用者进行参阅对比\par
\subsection{封面}
参照学校提供的统一封面，内容包括：学院名称，题目，学生姓名、学号，专业班级，指导教师姓名、职称。所有内容要居中。本模版封面按照学校提供的封面所制，排版上相同，但在打印上略有差别，使用封面时请自行斟酌。
\subsection{目录}
目录独立成页，包括论文中全部章、节的标题及页码。设计图纸要有标号。
\subsubsection{理工科类}

采用的三级标题序次结构有以下一种：\par
序次：1、1.1、1.1.1……\par
目录内容采用宋体五号。\par
注意：本模版目录格式按照为理工科要求所设，若需使用文科目录格式请自行修改。
\subsubsection{文科类}
目前采用的三级标题序次结构有以下两种：\par
第一种序次：第一章、第一节、一……\par
第二种序次：一、（一）、1……\par
目录采用宋体加粗四号居中，目录内容采用宋体五号，数字、字母采用Time New Roman五号。
\subsubsection{题目}
黑体三号居中，不宜超过20个汉字。
\subsubsection{摘要、关键词}
见模版摘要部分要求
\subsubsection{论文文体}
文本主体一般包括引言（或称前言、序言等）、正文与结论三部分。引言宣示课题的“来龙”，应说明本课题的意义、目的、主要研究内容、范围及应解决的问题。也可以不用引言，直接从第一章写起。正文是毕业论文（设计）的核心。结论部分可以用“结束语”、“结语”等标题来表示，也可以不用标题表示。\par



\section{图片插入}
\section{表格绘制}
\subsection{excel2latex插件}
https://ctan.org/tex-archive/support/excel2latex/
\subsection{格式调整}


\section{代码插入 }




% % 引用参考文献
% \clearpage
% \bibliographystyle{gbt7714-author-year}
% \bibliography{ref}
% \Reference{
% }
\Thanks
{
  岁月匆匆，大学时光如白驹过隙，我即将踏入毕业设计的最后阶段，心潮澎湃，回首往昔，感慨万千。在此，我要向所有曾经帮助、陪伴过我的人们深深地致以诚挚的谢意。\par
  首先，我要感谢母校，是她为我搭建了知识的殿堂，为我提供了探索未来的舞台。五年时光，荟萃了她对我的呵护与培育，我将永怀感恩之心。\par
  感谢建筑学院的恩师们，你们的悉心教导和引领，让我在学术之路上不断探索，不断超越。你们是我学习的灯塔，照亮了前行的道路。
}



\end{document}

